% Created 2024-06-16 Sun 14:04
% Intended LaTeX compiler: pdflatex
\documentclass[11pt]{article}
\usepackage[utf8]{inputenc}
\usepackage[T1]{fontenc}
\usepackage{graphicx}
\usepackage{longtable}
\usepackage{wrapfig}
\usepackage{rotating}
\usepackage[normalem]{ulem}
\usepackage{amsmath}
\usepackage{amssymb}
\usepackage{capt-of}
\usepackage{hyperref}
\author{cerebus2}
\date{\today}
\title{Certificate Manager Enrollment Over Secure Transport}
\hypersetup{
 pdfauthor={cerebus2},
 pdftitle={Certificate Manager Enrollment Over Secure Transport},
 pdfkeywords={},
 pdfsubject={},
 pdfcreator={Emacs 29.3 (Org mode 9.6.15)}, 
 pdflang={English}}
\begin{document}

\maketitle
\tableofcontents

\section{About}
\label{sec:org210998c}
\texttt{estissuer} is a \href{https://cert-manager.io/}{cert-manager} \href{https://cert-manager.io/docs/configuration/external/}{External certificate issuer} that implements a subset of \href{https://tools.ietf.org/html/rfc7030}{RFC 7030 Enrollment over Secure Transport} (EST).  It provides custom resources for EST services and a custom controller that reconciles certificate requests.

\texttt{est-operator} is written in Python using the \href{https://kopf.readthedocs.io/en/stable/}{kopf framework} and \href{https://lightkube.readthedocs.io/en/latest/}{lightkube}.

This operator covers only a portion of RFC 7030; see \href{docs/RFC7030.org}{this document} for an overview of what's implemented, planned, and not planned.

See \href{docs/DESIGN.org}{this file} for an overview of this operator's design.
\section{Getting Started}
\label{sec:org1faa682}
\subsection{Prerequisites}
\label{sec:orgdf6bf16}
\subsection{Building}
\label{sec:orgfad95ca}
\subsection{Installing}
\label{sec:orgfbb7fb3}
See \ref{sec:orgfbb80ee}, below.  Also see \href{docs/DESIGN.org}{the design document} for a deployment description.
\section{Testing}
\label{sec:orgc6b7dc8}
\subsection{Style}
\label{sec:org1760767}
\subsection{Component}
\label{sec:orgd7222ec}
\subsection{End-to-end}
\label{sec:orgd39ca9f}
\section{Deployment}
\label{sec:orgfbb80ee}
\subsection{In cluster}
\label{sec:org57b8ff8}
TBD.
\subsection{Out of cluster}
\label{sec:orgadf18ec}
\section{Usage}
\label{sec:orgb76633e}
\begin{itemize}
\item EstIssuer and EstClusterIssuer secrets must be passwords and labeled "est.mitre.org/registration=true".
\item Certificate secrets must be labeled "est.mitre.org/certificate=true".
\begin{itemize}
\item RFC 7030 renewals require mTLS authentication even when initial renewal uses a password regitration credential, so this label allows the controller to track certificate private keys.
\end{itemize}
\end{itemize}

\section{Federal Information Processing Standards (FIPS) 140-2 certified mode}
\label{sec:org99c2302}
This controller can be built to satisfy FIPS 140-2 certification requirements.  FIPS 140-2 certification depends on the module and the platform.  A platform with Python dynamically linked to OpenSSL \emph{and} links OpenSSL to a validated version of the OpenSSL FIPS Object Module can demonstrate FIPS 140-2 certification.  Consult the \href{https://csrc.nist.gov/projects/cryptographic-module-validation-program/validated-modules/search?SearchMode=Basic\&ModuleName=OpenSSL\&CertificateStatus=Active\&ValidationYear=0}{NIST Cryptographic Module Validation Program list of OpenSSL modules} for combinations of OpenSSL and platform that are certified.

\textbf{Note}: OpenSSL 1.0.2 is no longer supported by the OpenSSL project, but OpenSSL FIPS Object Module 2.0 is not compatible with OpenSSL 1.1.  A new version is \href{https://www.openssl.org/blog/blog/2018/09/25/fips/}{not available nor certified yet.}

Other FIPS certified modules are available (e.g., GnuTLS and Network Security Services) but the work to integrate these modules is extensive.  We welcome merge requests on this topic.

\subsection{Building a FIPS version}
\label{sec:org7f2ebc0}

\subsection{Building a FIPS container image}
\label{sec:org2e58573}

\subsection{Running a FIPS version out of cluster}
\label{sec:org33ac29c}

\section{Built with}
\label{sec:orgbbdf78a}

\section{Contibuting}
\label{sec:orge1205d5}
TBD
\section{Versioning}
\label{sec:orge329c77}
TBD
\section{Authors}
\label{sec:org29dba9c}
\begin{itemize}
\item Timothy J. Miller - \emph{initial work} - \href{https://github.com/Cerebus}{Cerebus}
\end{itemize}
\section{License}
\label{sec:orga6ffb93}
\href{./LICENSE}{Apache License v2.0}
\section{Acknowledgements}
\label{sec:org11c9880}
\begin{itemize}
\item \href{https://www.mitre.org}{The MITRE Corporation} for allowing this work to be open sourced.
\item \href{https://cert-manager.io}{Jetstack and all the cert-manager contributors} for the framework that makes this extension possible.
\item \href{https://github.com/nolar}{Sergey Vasilyev} for kopf, so I don't have to learn Go yet.
\end{itemize}
\end{document}
